\documentclass[11pt]{article}

\usepackage[margin=1in]{geometry}
\usepackage{amsmath,amssymb,amsthm}
\usepackage{booktabs}
\usepackage{microtype}

\newtheorem{theorem}{Theorem}
\newtheorem{corollary}{Corollary}
\newtheorem{remark}{Remark}

\newcommand{\bthree}{b_3}
\newcommand{\flux}{F}
\newcommand{\hh}{H}

\title{The flux landscape has exactly one exact structure: the D3 tadpole}
\author{Rob Mitchell}
\date{}

\begin{document}
\maketitle

\begin{abstract}
We prove that the D3 tadpole $N_{\rm flux}=\flux^\top J \hh$ is a \emph{complete}
invariant of the type IIB flux under the full duality group $\mathrm{Sp}(\bthree)
\times \mathrm{SL}(2,\mathbb{Z})$: two fluxes are dual-equivalent if and only if
they have the same tadpole. The proof is two lemmas. First (Mitchell's Uniqueness
Theorem, upgraded by the First Fundamental Theorem of invariant theory, Weyl 1939),
the ring of duality-invariant polynomials in the flux pair is generated by the
tadpole alone --- there is no invariant of any degree other than $N_{\rm flux}$ and
its powers. Second, the duality group acts transitively on each non-zero tadpole
level set (the orbit tangent space has rank $2\bthree-1$, equal to the level set's
codimension-$1$ dimension). Together they make $N_{\rm flux}$ a complete invariant.
Consequence --- a \emph{no-fingerprint no-go}: no exact, frame-independent quantity
can distinguish two vacua that share a tadpole, and since the tadpole is bounded by
the orientifold charge ($N_{\rm flux}\le L$), a landscape search can only
\emph{eliminate} (by the constraint), never \emph{select}. The exact content of a
theory of everything is thereby constrained to two classes --- consistency
constraints and topological labels --- and nothing else. The statement is
falsifiable: a single moduli-independent, duality-invariant flux quantity that is
not a function of $F^\top J H$ breaks it. All claims are verified symbolically
(sympy, zero tolerance) and reproduced by \texttt{scripts/validate\_all\_theorems.py}.
\end{abstract}

\section{Introduction}

The constraint-first program measured, on $174$ supersymmetric IIB flux vacua, that
the low-energy observables of a flux vacuum form a mostly-unstructured cloud: the
string coupling, the gauge couplings, the mass-spectrum shape, the axions, the
dilaton and the cosmological constant are statistically uniform in flux, while
exactly one quantity --- the D3 tadpole --- is exactly structured
\cite{constraint-first}. That was a measurement. This paper upgrades its central
claim --- \emph{``the landscape's only exact structure is what it forbids''} --- from
an empirical pattern to a theorem, by identifying the flux pair as the $k=2$ case of
a classical result in invariant theory.

\section{Setup}

Let $X$ be a Calabi--Yau threefold with middle cohomology $H^3(X,\mathbb{Z})\simeq
\mathbb{Z}^{\bthree}$, and let the two NS--NS/RR three-form fluxes be
$\flux,\hh\in H^3(X,\mathbb{Z})$. The intersection form is the antisymmetric
$J=\bigl(\begin{smallmatrix}0&I\\-I&0\end{smallmatrix}\bigr)$ in a symplectic basis,
and the D3 tadpole is the symplectic pairing
\begin{equation}
N_{\rm flux} \;=\; \flux^\top J \hh \;=\; \sum_{i=0}^{\bthree/2-1}\bigl(F_i H_{i+\bthree/2}-F_{i+\bthree/2}H_i\bigr).
\end{equation}
The duality group is $\mathrm{Sp}(\bthree,\mathbb{Z})$ (the monodromy group acting
on $H^3$, hence on $\flux$ and $\hh$ separately) together with S-duality
$\mathrm{SL}(2,\mathbb{Z})$ acting as $(\flux,\hh)\mapsto(a\flux+b\hh,c\flux+d\hh)$,
$ad-bc=1$. The tadpole is invariant under both: under $\mathrm{Sp}(\bthree)$ it is
the symplectic pairing, and under $\mathrm{SL}(2,\mathbb{Z})$ one computes
$N_{\rm flux}\mapsto(ad-bc)N_{\rm flux}$ because $\flux^\top J\flux=\hh^\top J\hh=0$
and $\hh^\top J\flux=-\flux^\top J\hh$.

\section{Mitchell's Uniqueness Theorem}

\begin{theorem}[Mitchell]
The D3 tadpole $\flux^\top J\hh$ is the unique non-trivial bilinear form in the flux
$(\flux,\hh)$ that is invariant under $\mathrm{Sp}(\bthree,\mathbb{Z})$, for any
$\bthree$.
\end{theorem}

The candidates $\flux^\top\flux$, $\hh^\top\hh$, $\flux^\top\hh$ are gauge-dependent
(they change under a symplectic transformation), while
$\flux^\top J\flux=\hh^\top J\hh=0$ identically by antisymmetry. The remaining
pairing is $\flux^\top J\hh$. This is verified symbolically ($6/6$ checks,
\texttt{scripts/validate\_mitchell\_uniqueness.py}).

\section{The Constraint-First Theorem: the full invariant ring}

Mitchell's theorem is the degree-$2$ slice. The whole curve is classical.

\begin{theorem}[Constraint-First, = First Fundamental Theorem, Weyl 1939]
Every $\mathrm{Sp}(\bthree)$-invariant polynomial in the flux pair $(\flux,\hh)$ is
a polynomial in the tadpole:
\begin{equation}
\mathbb{C}\bigl[(\mathbb{C}^{\bthree})^{2}\bigr]^{\mathrm{Sp}(\bthree)}
\;=\; \mathbb{C}\bigl[\flux^\top J \hh\bigr].
\end{equation}
\end{theorem}

This is Weyl's First Fundamental Theorem of invariant theory for the symplectic
group specialized to $k=2$ copies of the defining representation; the flux pair
$(\flux,\hh)\in\mathbb{C}^{\bthree}\oplus\mathbb{C}^{\bthree}$ is precisely that
$k=2$ case. It says the tadpole \emph{generates} the invariant ring: the only
invariants are $1,\ \flux^\top J\hh,\ (\flux^\top J\hh)^2,\ \dots$ We verify the
completeness direction computationally at every degree $d=1,\dots,5$ by annihilating
a torus-reduced ansatz with the full $\mathfrak{sp}(\bthree)$ algebra and computing
the nullity exactly; the invariant-space dimensions are
\begin{equation}
\dim = 0,\ 1,\ 0,\ 1,\ 0 \qquad (d=1,2,3,4,5),
\end{equation}
the signature of a ring generated by a single degree-$2$ element
(\texttt{scripts/validate\_constraint\_first.py}).

\begin{remark}
The honest attribution is that this is \emph{not new mathematics}: it is Weyl's
$1939$ theorem. The contribution here is the \emph{identification} --- that the flux
pair is the $k=2$ FFT case, so the ``only exact structure is the constraint'' is a
corollary of classical invariant theory rather than an empirical conjecture.
\end{remark}

\section{The No-Fingerprint No-Go Theorem}

\begin{theorem}[No-Fingerprint No-Go]
For $N_{\rm flux}\neq 0$, two flux pairs are dual-equivalent under
$\mathrm{Sp}(\bthree)$ if and only if they have the same tadpole. Equivalently,
$N_{\rm flux}$ is a \emph{complete} invariant.
\end{theorem}

\begin{proof}[Sketch]
Necessity is the Constraint-First Theorem: an invariant is a function of
$N_{\rm flux}$. Sufficiency is transitivity. The orbit tangent space at a pair is the
image of the map $\mathfrak{sp}(\bthree)\to\mathbb{R}^{2\bthree}$, $X\mapsto
(X\flux,X\hh)$, whose rank is $2\bthree-1$; this equals the codimension-$1$
dimension of the level set $\{\flux^\top J\hh=N\}$. The orbit is open in the level
set (equal tangent spaces), closed (an algebraic group orbit), and the level set is
connected for $N\neq 0$, so the orbit \emph{is} the level set.
\end{proof}

The rank is computed exactly (integer rref) at three generic pairs and at the
canonical pair $(e_0,e_3)$ itself, giving $2\bthree-1=11$ for $\bthree=6$
(\texttt{scripts/validate\_no\_fingerprint.py}).

\begin{corollary}[Search can only eliminate, never select]
No exact, frame-independent quantity can distinguish two vacua that share a tadpole.
The only exact selector is $N_{\rm flux}$ itself, and it is \emph{bounded}
($N_{\rm flux}\le L$ by the orientifold charge). Therefore a landscape search can
only \emph{eliminate} by the constraint --- it can never \emph{select} a special
vacuum, because there is no invariant left to select with.
\end{corollary}

\section{Full-duality completion}

\begin{theorem}
The invariant ring of the \emph{full} duality group $\mathrm{Sp}(\bthree)\times
\mathrm{SL}(2,\mathbb{Z})$ is also $\mathbb{C}[\flux^\top J\hh]$, so the no-go holds
for the complete duality group, not just the monodromy subgroup.
\end{theorem}

By the Constraint-First Theorem the $\mathrm{Sp}(\bthree)$-invariants are
$\mathbb{C}[N_{\rm flux}]$; the generator $N_{\rm flux}$ is $\mathrm{SL}(2,\mathbb{Z})$-invariant
(verified symbolically for the generators $S:(\flux,\hh)\mapsto(-\hh,\flux)$ and
$T:(\flux,\hh)\mapsto(\flux,\flux+\hh)$, and for a general $\mathrm{SL}(2)$ element
$N_{\rm flux}\mapsto(ad-bc)N_{\rm flux}$), so $\mathrm{SL}(2,\mathbb{Z})$-invariance
removes nothing. The full invariant ring is exactly $\mathbb{C}[N_{\rm flux}]$
(\texttt{scripts/validate\_full\_duality.py}).

\section{The Complete Exact Structure Theorem}

\begin{theorem}[Complete Exact Structure]
For a $\mathrm{CY}_3$ orientifold flux compactification, every exact
(moduli-independent, duality-invariant) observable is either (i) the D3 tadpole
$N_{\rm flux}$ --- the unique flux observable --- or (ii) a topological invariant of
the Calabi--Yau ($h^{1,1},h^{2,1},c_2,\kappa$). There is nothing else.
\end{theorem}

In particular $W_0$, $g_s$, the masses, the saxion values, the gauge couplings and
the cosmological constant are \emph{not} exact observables: they are moduli- or
flux-dependent, and duality-frame-dependent, hence exactly the quantities a theory
of everything does \emph{not} get to predict.

\begin{corollary}[Hard ceiling on the verification deficit]
The flux sector contributes exactly \emph{one} observable to the Mitchell
Verification Deficit $m$. By the no-go there is provably no second flux observable
to find, so $m$ can never grow from the flux sector; future growth must come from
the other sectors (D-brane charges, the Standard Model), which is the $m=0$ frontier.
\end{corollary}

\begin{remark}[Falsifiability]
The whole tower stakes itself on one checkable claim: \emph{no moduli-independent,
duality-invariant flux quantity exists that is not a function of $\flux^\top J\hh$.}
A single counterexample falsifies it. That is what makes this a theorem rather than
a slogan.
\end{remark}

\section{Conclusion}

The constraint-first conjecture is now a theorem. The flux sector's only exact
structure is the D3 tadpole --- a \emph{constraint}, not a selector --- and the
theory of everything predicts nothing exactly except what it forbids (consistency
constraints) and what it is handed (topology). Search is provably dead; the only
navigable move is elimination by verified constraint.

\paragraph{Scope, named.} The invariant-ring claim is proven for \emph{polynomial}
invariants (the FFT); smooth non-polynomial invariants are not covered by the
symbolic proofs, though nothing suggests they exist. Global transitivity rests on
two classical facts (the orbit is closed; the $N\neq0$ level set is connected), stated
as such; their machine-checked (Lean) formalization is deferred to hardware. The
topology list $h^{1,1},h^{2,1},c_2,\kappa$ is the standard complete topological data
of a $\mathrm{CY}_3$, taken as established mathematics.

\begin{thebibliography}{9}
\bibitem{constraint-first} R.~Mitchell, ``Constraint-first string theory,'' companion paper (this series).
\bibitem{weyl} H.~Weyl, \emph{The Classical Groups: Their Invariants and Representations}, Princeton University Press (1939).
\bibitem{denef} F.~Denef and M.~R.~Douglas, ``Distributions of flux vacua,'' JHEP \textbf{0405} (2004) 072.
\bibitem{jaxvacua} A.~Dubey, S.~Krippendorf, A.~Schachner, ``JAXVacua'' (arXiv:2306.06160).
\bibitem{fulton} W.~Fulton and J.~Harris, \emph{Representation Theory: A First Course}, Springer (1991).
\end{thebibliography}

\end{document}
